\chapter{The 3D camera}

\section{It is five numbers, and none of them are magic}

\begin{lstlisting}
typedef struct Camera3D {
    Vector3 position;     // where the eye is
    Vector3 target;       // the world point it is looking at
    Vector3 up;           // which way is "up" for the camera
    float fovy;           // vertical field of view, in DEGREES
    int projection;       // CAMERA_PERSPECTIVE or CAMERA_ORTHOGRAPHIC
} Camera3D;
\end{lstlisting}

\begin{center}
\begin{tikzpicture}[scale=1.25,every node/.style={font=\sffamily\footnotesize}]
  \fill[rlight,rounded corners=3pt] (-1.2,-1.6) rectangle (7.6,2.4);

  % eye
  \fill[rink] (0,0) circle (2.4pt);
  \node[left,text=rink] at (-0.12,0) {\texttt{position}};

  % target
  \fill[rorange] (5.4,0.6) circle (2.6pt);
  \node[right,text=rorange] at (5.6,0.6) {\texttt{target}};

  % view direction
  \draw[->,rgrey,dashed] (0,0) -- (5.3,0.59);
  \node[below,text=rgrey] at (2.7,0.18) {direction = normalize(target - position)};

  % frustum
  \draw[rblue,line width=0.9pt] (0,0) -- (6.4,1.9);
  \draw[rblue,line width=0.9pt] (0,0) -- (6.4,-0.75);
  \draw[rblue,->] (1.6,0.48) arc (16.5:-8:1.7);
  \node[right,text=rblue] at (1.85,0.05) {\texttt{fovy}};

  % up
  \draw[->,rgreen,line width=1.2pt] (0,0) -- (0,1.5);
  \node[above,text=rgreen] at (0,1.5) {\texttt{up} = (0,1,0)};
\end{tikzpicture}
\end{center}

Notice what is \emph{not} in that struct: there is no field for the direction the
camera is facing, and no rotation angles. The direction is worked out from the
two points:

\[ \text{direction} = \text{normalize}(\text{target} - \text{position}) \]

Which gives you the two operations you will use forever:

\begin{itemize}[leftmargin=*,itemsep=2pt]
  \item \textbf{To turn the camera}: move the \lstinline|target|, leave
        \lstinline|position| alone.
  \item \textbf{To move the camera}: move \lstinline|position| \emph{and}
        \lstinline|target| by the same amount, so it does not turn while it
        travels.
\end{itemize}

\gotcha{Move only \lstinline|position| and the camera swings round to keep
staring at the old target, like a security camera tracking you. That is
occasionally what you want and usually a bug.}

\section{fovy, or choosing a lens}

\lstinline|fovy| is the vertical angle the camera can see, in degrees. It behaves
exactly like a camera lens:

\begin{center}
\begin{tabular}{@{}ll@{}}
\toprule
\textbf{fovy} & \textbf{Feels like} \\
\midrule
20--30 & telephoto: flat, compressed, distant \\
45 & neutral, good for inspecting an object \\
60--75 & the normal first person game range \\
90+ & fisheye: fast, dizzying, very wide \\
\bottomrule
\end{tabular}
\end{center}

Doom-like shooters tend to sit around 70 to 90, which is why the mini Doom in
chapter 14 uses 75.

\section{Perspective versus orthographic}

\lstinline|CAMERA_PERSPECTIVE| is what your eyes do: distant things look smaller,
parallel lines converge.

\lstinline|CAMERA_ORTHOGRAPHIC| removes that. Distance stops affecting size and
parallel lines stay parallel. This is the look of isometric strategy games, of
CAD, and of 2D games faking depth.

\gotcha{When you switch to orthographic, \lstinline|fovy| stops meaning degrees
and starts meaning \emph{how many world units fit vertically on screen}. Leave it
at 45 and your world is 45 metres tall in the window --- everything is a speck.
Set something like 10 and it looks sane again.}

\section{Begin and End, again}

\begin{lstlisting}
BeginMode3D(camera);
    DrawGrid(20, 1.0f);
    DrawCube((Vector3){ 0, 1, 0 }, 2, 2, 2, MAROON);
EndMode3D();

// back to 2D screen pixels: the HUD
DrawText("HEALTH 100", 20, 20, 20, RAYWHITE);
DrawFPS(10, 10);
\end{lstlisting}

Exactly the same contract as \lstinline|BeginMode2D| in chapter 8: inside is
world space, outside is screen pixels. Your crosshair, health bar, ammo counter
and minimap all go outside, and that is why they stay nailed to the window while
the world moves.

There is one extra thing \lstinline|BeginMode3D| does that the 2D version does
not: it enables the \textbf{depth buffer}. In 2D the draw order decided what was
on top; in 3D the GPU compares distances per pixel, so a cube behind another cube
is hidden correctly no matter which order you drew them in. You get that for
free.

\gotcha{Transparency is the exception. The depth buffer works per pixel, and a
half transparent surface writes depth as if it were solid, which can hide things
behind it. If you draw transparent objects, draw them after the opaque ones,
furthest first. You will not hit this for a while, but when a glass wall makes
things vanish, this is why.}

\section{The ready-made cameras}

Before writing your own, try raylib's:

\begin{lstlisting}
UpdateCamera(&camera, CAMERA_FREE);           // fly around: WASD, mouse, wheel
UpdateCamera(&camera, CAMERA_ORBITAL);        // spins around the target by itself
UpdateCamera(&camera, CAMERA_FIRST_PERSON);   // classic FPS
UpdateCamera(&camera, CAMERA_THIRD_PERSON);   // follows the target from behind
\end{lstlisting}

One line, in the update section, and the camera works. For first and third
person, call \lstinline|DisableCursor()| first so the mouse is captured.

They are perfect for looking at a model or prototyping a scene, and they run out
of road the moment you want collisions, crouching, recoil, a vehicle, or a camera
that does anything specific to your game. That is the subject of the next
chapter.

\tryit{Lesson 2 of the 3D course lets you move each of the five camera fields
with the keyboard and watch the numbers; lesson 3 cycles the four built-in modes
with TAB.}{./course3d/build.sh 2 \quad and \quad ./course3d/build.sh 3}
